\section{Smart Contracts}\label{ch:contracts}

This chapter discharges constitution clauses C-6.1--C-6.5 (smart contracts as converter, on-record inputs, purity and statelessness, and every right or obligation as a unit --- C-6.5 jointly with Chapters~\ref{ch:collateral} and~\ref{ch:cdm}).

The ledger admits well-formed transactions and records them; it performs no
product-specific logic of its own. What a dividend, an expiry, a barrier knock,
or a settlement print means for a particular instrument lives in that
instrument's \emph{smart contract}. An event arrives from the outside or is
emitted by the Event Monitor; the ledger can do nothing with it until it is a
transaction; the smart contract is the one function that performs the
conversion. Since the ledger is a transaction processor and nothing more,
everything reaching it from outside must be turned into a transaction, and the
contract is the converter.

One signature states the entire interface.

\begin{lstlisting}
contract :: Event -> Ledger -> Transaction

-- what a contract is NOT (each of these would be a defect):
--   contract :: Event -> State Memory Transaction   -- it may not remember
--   contract :: Event -> IO Transaction             -- it may not reach out
--   contract :: Event -> Clock -> Transaction        -- it may not read time
\end{lstlisting}

Read the signature as a claim about what a contract may touch. Its only inputs
are the event and the ledger as it currently stands: the declared terms of the
instrument, the recorded observations, and the visible unit and position state
at the moment of firing. There is no clock among the inputs --- the clock is
consulted once, by the Event Monitor of Chapter~4, and never again downstream
--- no private store, and no channel to the world. A smart contract is
therefore a pure, deterministic function of its inputs, and, because none of
those inputs is private to it, a stateless one. Each contract is exactly this: a
stateless, machine-readable transcription of the legal terms and conditions of
the instrument it governs. Purity is what makes economic verification possible
at all: to check a recorded transaction, re-run the same contract on the same
recorded event and compare (Chapter~15). A function that read a clock, a store,
or the network could not be re-run to the same answer, and the guarantee would
be lost.

Two principles govern every contract, and the rest of this chapter is their
working-out.

\begin{principle}[Everything a contract consumes is on the record]
Every input a smart contract reads --- the instrument's terms, the observations
it depends on, the positions it acts upon --- is present on the record at the
moment of firing. A contract reads nothing that is not recorded, so its output
is a function of the record alone.
\end{principle}

\begin{principle}[A contract never waits and never remembers]
The waiting is done by the watch, which sits in the record, and by the Event
Monitor, which notices when the watch's condition is met; the contract is
invoked only after its event has fired, runs once, and returns. Whatever a
firing must carry forward to a later firing it writes into the record --- into
one of the three homes of Chapter~6 --- and the later firing reads it back. A
contract that had to remember would need private state, and private state is a
second store of the very facts the ledger already holds: the one thing this
architecture exists to abolish.
\end{principle}

\subsection{The dividend: three firings, one contract}

The dividend shows both principles in one instrument fired three times. Wallet
W-ALPHA holds 10{,}000 shares of the stock ACME.

\emph{The trigger.} On 2026-05-04 the issuer announces a cash dividend of
\USD{1.50} per share on ACME. The announcement is captured at the boundary and
recorded --- an external event, like a trade.

\emph{The contract that catches it.} The stock's smart contract fires on the
announcement, and fires again on each date the announcement creates. The three
firings are one contract, invoked three times on three different events; it
carries nothing between them.

\emph{The transactions it produces.}

\begin{center}
\begin{tabular}{clll}
\toprule
Firing & Date & Event & Transaction emitted \\
\midrule
1 & 2026-05-04 & announcement & arm the record-date and payment-date watches; \\
  &            &              & \emph{no move} \\
2 & 2026-05-15 & record date  & PositionState: W-ALPHA entitlement \\
  &            &              & $= 10{,}000 \times \USD{1.50} = \USD{15{,}000.00}$; \emph{no move} \\
3 & 2026-05-22 & payment date & move \USD{15{,}000.00}, issuer $\rightarrow$ W-ALPHA; \\
  &            &              & PositionState: entitlement discharged \\
\bottomrule
\end{tabular}
\end{center}

The first firing pays nothing; it registers the two watches the dividend
creates and writes them into the unit's state. The second reads the owned
positions as the ledger then stands and records each holder's entitlement --- a
moveless transaction: no cash changes hands, but the position state now carries
what each holder is owed. The third reads those recorded entitlements and emits
the cash moves against them. Nothing links the firings but the record: firing~2
does not remember firing~1, it reads the armed watch; firing~3 does not remember
firing~2, it reads the recorded entitlement.

\emph{What the door checks.} On firing~3 the Transaction Executor checks
conservation --- the \USD{15{,}000.00} debited from the issuer wallet equals the
amount credited to W-ALPHA, a paired-leg move that can neither create nor
destroy quantity --- and idempotence: the payment proposal carries an identifier
derived from its cause, so a retried or late firing commits once (Chapter~4).

\emph{Two dividends in flight.} A second dividend can be declared before the first has paid.
ACME declares again on 2026-06-09, payable 2026-06-30, while the 2026-05-22 payment is still
a week off. The record-date firing of the second dividend records a second entitlement for
W-ALPHA, and the position state now carries both at once. They do not collide: each
entitlement is a separate entry in PositionState, keyed by the cause-derived identifier of the
firing that wrote it (Chapter~\ref{ch:homes}), so the payment-date firing of the first
discharges the first and retires its entry, leaving the second live until its own payment
date. Two entitlements in flight is a routine Tuesday, and the position state holds one live
fact for each.

\emph{Where an agreement encumbers the holder} --- the shares pledged, the cash
trapped --- the payment routes through a conditional obligation unit instead of
paying directly; that refinement is Chapter~9's, and it changes what firing~3
emits, never that the contract is stateless and fired three times.

\emph{Three firings, three transactions, one stateless contract; between them
the contract remembers nothing --- the armed watches, the recorded entitlement,
and the positions are all on the record.}

\subsection{The trigger classes differ in what they read, never how they run}

The three firings of the dividend were reached three different ways. The
announcement arrived from outside. The record date and the payment date were
calendar conditions the Event Monitor noticed against the clock. A fourth way
remains: a predicate on a recorded observation --- a settlement price printing,
a close falling below a barrier. These are the trigger classes, and they
classify what a contract reads, never how it runs.

\begin{lstlisting}
data Trigger
  = External  Observation  -- captured at the boundary (a dividend announcement)
  | Calendar  Date         -- a date arriving (a record, payment, or expiry date)
  | Condition Predicate    -- a predicate on a recorded observation (a print, a touch)
-- every trigger, whatever its class, fires the same  contract  above.
\end{lstlisting}

The class decides which watch the Monitor evaluates and which inputs the firing
reads; the signature above is unchanged across all of them. A future exercises
two of the classes on one instrument.

\subsection{The future: settlement prints and the expiry unwind}

FUT-IDX is a cash-settled future on the index IDX, multiplier~10, settled to
market. W-ALPHA holds one contract, entered at 995 against the clearing house
wallet. Its terms declare a watch on each daily settlement print and a watch on
its expiry.

\emph{The trigger and the contract.} On each recorded settlement print the
Monitor emits and the future's contract fires. It reads two facts, each already
on the record: the day's settlement level, a UnitStatus fact identical for every
holder, and W-ALPHA's last applied level, a PositionState fact. From the two it
computes the day's variation margin and emits it. On the third print, which is
also the expiry, the same contract additionally unwinds the position --- in one
atomic transaction.

\begin{center}
\small
\begin{tabular}{clrll}
\toprule
Day & Print & Variation margin & Move & PositionState \\
\midrule
1 & 1{,}000.00 & $(1000-995)\times 10 = +50$   & \USD{50}, CCP $\rightarrow$ W-ALPHA  & level $\leftarrow 1{,}000.00$ \\
2 & \phantom{0}990.00 & $(990-1000)\times 10 = -100$ & \USD{100}, W-ALPHA $\rightarrow$ CCP & level $\leftarrow 990.00$ \\
3 & 1{,}005.00 & $(1005-990)\times 10 = +150$  & \USD{150}, CCP $\rightarrow$ W-ALPHA & position unwound \\
\bottomrule
\end{tabular}
\end{center}

The contract never remembers yesterday's level; it reads it from PositionState,
computes against today's print, and writes today's level back for tomorrow's
firing. Cumulative variation margin is $+50-100+150 = +100 = (1005-995)\times 10$
--- exactly the position's profit, arrived at one recorded firing at a time.

\emph{More than one lot.} The single-contract reading above stores one level, but a holder
can carry several lots entered on different days, and a new lot can trade between two prints.
So PositionState does not store a bare level; it stores the \emph{aggregate settled value}
$S = m\sum_i q_i\cdot\text{level}_i$ as one integer, for multiplier $m$ --- the value the
position's lots have already been marked to (Chapter~\ref{ch:homes}). The day's variation
margin is then $m\,\text{settle}_t\,Q - S$, where $Q=\sum_i q_i$ is the total quantity. For
one lot at level~$995$ this is $S = 10\times 995 = 9{,}950$ and day~1 gives
$10\times 1{,}000 - 9{,}950 = +50$, the table above. When a new lot of $q'$ trades at price
$p$ before the print, the firing folds it in at its trade price, $S \leftarrow S + m\,q'\,p$:
if W-ALPHA trades a second contract at $992$ before day~2's print, $S = 10{,}000 + 10\times
992 = 19{,}920$ and $Q=2$, so day~2's margin is $10\times 990\times 2 - 19{,}920 = -120$ ---
the old lot marked from yesterday's $1{,}000$ ($-100$) and the new lot from its own trade
price $992$ ($-20$). The whole position needs a single stored integer, not one entry per lot.
This is what a clearing house keeps, and it is why the futures fact rides inside the
position-state collection (Chapter~\ref{ch:homes}) as one aggregate entry rather than
accumulating a fact per lot.

\emph{What the door checks.} The margin move is conserved --- what W-ALPHA
receives the clearing house pays, and conversely --- and each firing multiplies
the multiplier by \emph{the recorded settlement level}, never a stale one, so no
firing books a quantity against a price from a different state (Chapter~14).

\emph{The daily print and the expiry unwind are the same contract fired on
different events; the settlement print and the dividend announcement differ in
what the contract reads, never in how it runs.}

\subsection{The one-touch: one firing, one transaction, and idempotence}

OT-1 (Cast, \S\ref{sec:cast}) declares two watches on OMEGA. The first is on
OMEGA's recorded official close against the barrier --- and because the barrier is observed at that close,
the watch is exactly the barrier edge's declared guard, so the term is
transcribed faithfully (Chapter~3's product-graph consistency). The second is
OT-1's corporate-action out-edge: because OT-1's barrier references OMEGA, its
terms watch OMEGA for a corporate action, and should OMEGA split two-for-one this
edge fires, appends a new ProductTerms version carrying the barrier
$80.00 \rightarrow 40.00$ under the term-adjustment operator
(Chapter~\ref{ch:marketdata}), and leaves OT-1 on \emph{live} --- a self-loop that
bumps the terms version, not the state (Chapter~3,
Section~\ref{sec:termsheet-graphs}). Without it the first post-split close near
$47.50$ would knock a barrier struck at $80.00$ against a halved stock, a phantom
the adjusted barrier $40.00$ refuses.

\emph{The trigger and the contract.} OMEGA closes at 79.00; the close is
recorded, the condition is met, and the Monitor emits the touch. The contract
fires and reads exactly two things from the record: the declared payout and who
holds the option. It returns one transaction --- a single move of
\USD{1{,}000{,}000} from W-BANK to W-ALPHA, together with one unit-state change,
OT-1 $\rightarrow$ \emph{triggered}.

\emph{What the door checks.} Conservation on the payout, and idempotence.
Idempotence is not something the contract defends: should the Monitor emit the
touch a second time --- barriers are often crossed on many later days --- the
next firing reads the unit and finds it already \emph{triggered}, so it proposes
nothing, and even a duplicate of the original proposal carries the same
cause-derived identifier and commits once at the door.

\emph{Idempotence is guaranteed not by the contract's vigilance but by the
recorded state and the cause-derived identifier; the touch contract is a pure
function of the record --- terms and holder in, one transaction out.}

\subsection{The variance swap: a contract never remembers}

VS-1 is a one-year variance swap on IDX with 252 scheduled fixings against the
same recorded official closes the future settles on. Each scheduled fixing is a
declared watch; on fixing~$k$ the Monitor emits and the contract fires. To
extend the realised variance the firing needs the previous close and the running
accrued statistic --- the fixings $1 \ldots k-1$. It does not hold them: the
previous close is a recorded observation the firing reads from the record --- a
home fact folded from its observation-recording transaction, never a live feed ---
and the accrued statistic is a UnitStatus fact that firing~$k-1$ wrote for exactly
this purpose (Chapter~6).
Firing~$k$ reads both, forms the return, adds its square, and writes the updated
accrued statistic back for firing~$k+1$. A firing 200 days into the swap's life
reconstructs its entire history from the record and holds nothing of its own
between firings.

\emph{Ordering.} The fold obeys the log's order, not the wire's. A fixing carries
two coordinates: its valid-time index~$k$ --- the fixing-sequence position it
belongs to --- and its transaction-time position, where its print was appended to
the log (Chapter~\ref{ch:picture}, Section~\ref{sec:bitemporal}). The accrual folds
in valid-time order, so replay is deterministic regardless of arrival order. A
duplicate print is inert, absorbed once under its cause-derived identifier; a late
or out-of-order print --- one whose transaction-time position trails its valid-time
index --- is not a no-op: it is folded at its observation position~$k$, and the
accrued-statistic tail $A_k, A_{k+1}, \ldots$ is recomputed from~$k$ forward. The
recomputed result is identical to timely arrival because the fold is over the
valid-time sequence, while the as-known-at accrual reported before the print
arrived is preserved beside it; any figure already posted against a superseded tail
is corrected by a compensating transaction, never left standing (Chapter~15).

\emph{Firing~$k$ finds fixings $1 \ldots k-1$ on the record; a contract never
remembers.}

\subsection{Faithfulness: every right and obligation is a unit}

Purity and statelessness say how a contract runs. One further demand says what it
must produce: the alignment with the legal contract runs deep enough that every
right and obligation the legal contract creates is represented in the ledger as
a unit. A subscription is booked as an exchange --- cash against the client's
redemption claim --- never as a transfer for nothing. A dividend not yet paid is
a recorded entitlement, seen above. A payout trapped by a collateral agreement
is a conditional obligation unit (Chapter~9). Because each such right and
obligation is a unit or a declared term of one, the contract that transcribes
the instrument has, at every firing, a home on the record for everything it must
carry --- which is why it can afford to remember nothing. Statelessness and
faithfulness are the same discipline seen from two sides: the contract keeps no
private state precisely because the record already holds, as units and their
declared terms, every fact its next firing will need.
